\documentclass[a4paper,UTF8]{ctexart}
\usepackage{settings}
\usepackage{hyperref}

\begin{document}

\sectionTitleI

% --------- %

\begin{listWithRemark}
\begin{exampleSentence}{1}
    {Elle est arrivée à Paris il y a deux jours pour passer deux semaines chez sa tante, la famille Moreau.}
    {她两天前抵达巴黎，为的是在她姑妈莫罗家住两个星期。}
    {}&{\hei 句子结构}&{主语 Elle + 谓语 est arrivée（复合过去时）+ 地点状语 à Paris + 时间状语 il y a deux jours + 目的状语 pour + inf.}\\
    {}&{\hei 语法结构}&{arriver 在复合过去时中用 être 作助动词，过去分词与主语阴性单数配合，故写作 arrivée。}\\
    {}&{\hei 考点提示}&{注意 il y a 表示"多少时间以前"，与 depuis 的持续义不同。}
\end{exampleSentence}
&
\exampleSentenceRemark[il y a deux jours 表示"两天前"这个过去时间点，不表示从过去持续到现在。]
\\
\begin{exampleSentence}{2}
    {La famille Moreau habite dans un quartier calme, près du centre-ville.}
    {莫罗一家住在一个安静的街区，靠近市中心。}
    {}&{\hei 句子结构}&{主语 La famille Moreau + 谓语 habite + 地点状语 dans un quartier calme + 补充地点 près du centre-ville.}\\
    {}&{\hei 语法结构}&{près de 表示"靠近"，后面接名词时别漏掉 de。}
\end{exampleSentence}
&
\exampleSentenceRemark[如果你想让右栏更像讲义批注，可以把这一栏只保留关键词，不写成长句。]
\\
\begin{exampleSentence}{3}
    {Charlotte et sa cousine Patricia sont parties de bonne heure pour prendre le métro Ligne~8.}
    {夏洛特和她的表姐帕特里西亚一大早就出发了，为了乘坐8号地铁线。}
    {}&{\hei 句子结构}&{主语（两人）+ 谓语 sont parties（复合过去时）+ 时间状语 de bonne heure + 目的状语 pour + inf.}\\
    {}&{\hei 语法分析}&{\ding{172} partir 复合过去时助动词用 être，主语为两个女性 → parties（阴性复数）；\ding{173} de bonne heure 固定表达，意为"早早地"；\ding{174} pour + prendre 目的状语}\\
    {}&{\hei 考点提示}&{être 助动词动词（位移动词）的复合过去时；过去分词性数配合规则}\\
\end{exampleSentence}
&
\exampleSentenceRemark[主语 Charlotte et Patricia 为两人，过去分词需用复数 parties，且两人均为女性，需加阴性 -s。]
\\
\begin{exampleSentence}{4}
    {Vingt minutes après, elles sont descendues à Créteil.}
    {二十分钟后，她们在克雷泰伊站下了车。}
    {}&{\hei 句子结构}&{时间状语 Vingt minutes après + 主语 elles + 谓语 sont descendues + 地点状语 à Créteil}\\
    {}&{\hei 语法分析}&{descendre 此处为不及物用法（下车），助动词用 être；à Créteil 表"在克雷泰伊（站）"}\\
    {}&{\hei 考点提示}&{双助动词动词（descendre/monter/sortir/rentrer 等）的助动词选择}\\
\end{exampleSentence}
&
\exampleSentenceRemark[descendre 可及物也可不及物：及物（搬下某物）时用 avoir，不及物（下来/下车）时用 être；本句为不及物用法。]
\\
\begin{exampleSentence}{5}
    {La pauvre fille n'est jamais entrée dans un si grand magasin, elle est bien étonnée.}
    {这个可怜的女孩从未走进过这么大的商场，她非常惊讶。}
    {}&{\hei 句子结构}&{主语 + ne...jamais 否定复合过去时 + dans un si grand magasin（地点）；第二分句：主语 + être + adj.}\\
    {}&{\hei 语法分析}&{n'est jamais entrée：否定复合过去时，否定词 ne...jamais 夹住助动词 est；un si grand magasin：si 此处为程度副词"如此/这么"，语序为 un/une + si + adj. + n.；bien 修饰 étonnée 表程度"非常"}\\
    {}&{\hei 考点提示}&{复合过去时否定结构；程度副词 si 的位置}\\
\end{exampleSentence}
&
\exampleSentenceRemark[复合过去时否定：ne 在助动词前，jamais 在助动词后，过去分词在最后；si 与 aussi 的区别：否定句程度副词用 aussi→si，此处为肯定句强调。]
\\
\begin{exampleSentence}{6}
    {C'est la première fois qu'elle voit un hypermarché.}
    {这是她第一次见到大型超市。}
    {}&{\hei 句子结构}&{C'est + la première fois + que 引导的关系从句（现在时）}\\
    {}&{\hei 语法分析}&{c'est la première/dernière fois que + 现在时：固定句型，从句使用现在时而非过去时，以说话时刻为参照}\\
    {}&{\hei 考点提示}&{c'est la première fois que 句型；关系从句时态选择}\\
\end{exampleSentence}
&
\exampleSentenceRemark[许多学习者误用过去时 voyait 或 a vu，正确用法是现在时 voit。]
\\
\begin{exampleSentence}{7}
    {À la campagne, Charlotte et sa mère doivent aller dans plusieurs petites boutiques pour faire leurs achats quotidiens.}
    {在农村，夏洛特和她的母亲必须去几家小店铺才能完成日常采购。}
    {}&{\hei 句子结构}&{地点状语 à la campagne + 主语 + 情态动词 doivent + inf. + 地点 + 目的状语}\\
    {}&{\hei 语法分析}&{devoir + inf. 情态动词"必须"；pour faire 目的状语；achats quotidiens 日常采购，achats 为复数}\\
    {}&{\hei 考点提示}&{情态动词 devoir 的用法}\\
\end{exampleSentence}
&
\exampleSentenceRemark[]
\\
\begin{exampleSentence}{8}
    {Il y a beaucoup de grands magasins comme ça dans la région parisienne, lui explique Patricia.}
    {巴黎地区有很多这样的大商场，帕特里西亚向她解释道。}
    {}&{\hei 句子结构}&{引用句 + 倒装引语动词：...lui explique Patricia（主谓倒装）}\\
    {}&{\hei 语法分析}&{lui explique Patricia：引语后置的倒装结构；lui 为间接宾语代词，指夏洛特；expliquer qch à qn → lui expliquer qch}\\
    {}&{\hei 考点提示}&{引语倒装；间接宾语代词位置}\\
\end{exampleSentence}
&
\exampleSentenceRemark[间接宾语代词 lui 放在动词前；引语倒装时主语后置。]
\\
\begin{exampleSentence}{9}
    {Oui, on y vend de tout, et tout est moins cher.}
    {是的，那里什么都卖，而且什么都更便宜。}
    {}&{\hei 句子结构}&{on y vend de tout；tout est moins cher（比较级）}\\
    {}&{\hei 语法分析}&{y 代替前文提到的地点（商场）；vendre de tout = 卖各种东西（de 为泛指用法）；moins cher 比较级，此处为绝对比较，无明确比较对象}\\
    {}&{\hei 考点提示}&{代词 y 的用法；比较级 moins + adj.}\\
\end{exampleSentence}
&
\exampleSentenceRemark[]
\\
\begin{exampleSentence}{10}
    {Ensuite, elles vont choisir des fruits et des légumes sans oublier des gâteaux.}
    {然后，她们去挑选水果和蔬菜，也没忘记蛋糕。}
    {}&{\hei 句子结构}&{时间副词 Ensuite + 主语 + aller + inf.（近将来）+ sans + inf.（方式状语）}\\
    {}&{\hei 语法分析}&{aller + inf. 近将来时，表"即将去做某事"；sans oublier + n. 固定表达"不忘……"；sans + inf. 表"没有……地"}\\
    {}&{\hei 考点提示}&{aller + inf. 近将来时；sans + inf. 结构}\\
\end{exampleSentence}
&
\exampleSentenceRemark[]
\\
\begin{exampleSentence}{11}
    {Chez nous, quand on a besoin de quelque chose, on va dans les petites boutiques.}
    {在我们那里，每当需要什么东西，就去小店铺买。}
    {}&{\hei 句子结构}&{地点状语 chez nous + 时间从句 quand + 主句 on va dans...}\\
    {}&{\hei 语法分析}&{chez nous 表"在我们那个地方（家乡）"，不是"在我们家"；avoir besoin de qch 固定搭配；on 泛指"人们"，口语中极常用；quand + 现在时表条件/习惯}\\
    {}&{\hei 考点提示}&{avoir besoin de 搭配；on 的泛指用法；chez 后接人的用法}\\
\end{exampleSentence}
&
\exampleSentenceRemark[]

\end{listWithRemark}

\newpage
\sectionTitleII

以下搭配均来自课文：

\begin{collocationTable}
{\hei 常用搭配} & {\hei 释义} & {\hei 用法说明} \\
{habiter dans qch} & {住在…} & {Elle est arrivée à Paris il y a deux jours.} \\
\SetCell[c=3]{l} \iiex{Elle est arrivée à Paris il y a deux jours pour passer deux semaines chez sa tante, la famille Moreau.} \\
{arriver à + 地点} & {抵达某地} & {arrivée à Paris} \\
\SetCell[c=3]{l} \iiex{Elle est arrivée à Paris il y a deux jours pour passer deux semaines chez sa tante} \\
{partir de bonne heure} & {一大早出发} & {sont parties de bonne heure} \\
\SetCell[c=3]{l} \iiex{Il faut partir de bonne heure.} \\
{prendre le métro} & {乘地铁} & {pour prendre le métro Ligne 8} \\
\SetCell[c=3]{l} \iiex{Je prends le bus tous les jours.} \\
{descendre à + 站名} & {在某站下车} & {descendues à Créteil} \\
\SetCell[c=3]{l} \iiex{Descendez à la prochaine station.} \\
{entrer dans qch} & {走进……} & {Il entre dans le bureau} \\
\SetCell[c=3]{l} \iiex{Je n'est jamais entrée dans un grand magasin.} \\
{prendre un kilo de} & {取一公斤……} & {prennent un kilo de viande} \\
\SetCell[c=3]{l} \iiex{Prenez deux kilos de tomates.} \\
\end{collocationTable}

\newpage
\sectionTitleIII

\begin{listWithRemark}
\begin{vocabEntry}
    {digérer}{\posvt}{消化；（转）领会，理解}
    {}&{\hei 常用搭配}&{
        digérer un repas \hspace{6pt} 消化一顿饭 \\
        difficile à digérer \hspace{6pt} 难消化的 \\
        mal digérer une information \hspace{6pt} 难以理解某信息}\\
    {}&{\hei 变位信息}&{
        第三组 -érer 动词（按 espérer 模式变位）\\
        直陈式现在时：digière, digères, digère, digérons, digérez, digèrent；\\
        复合过去时：avoir digéré；\\
        简单将来时：digérerai, digéreras...；\\
        虚拟式现在时：digière, digères...}\\
\end{vocabEntry}
&
\begin{vocabRemark}
    \vocabMot{digestion}{\sposnf}{消化} \\
    \vocabMot{digestif}{\sposadj}{助消化的}
\end{vocabRemark}
\\
\begin{vocabEntry}{individuel}{\posadj}{个人的，个体的}
    {}&{\hei 常用搭配}&{
        plat individuel \hspace{6pt} 个人份菜肴 \\
        chambre individuelle \hspace{6pt} 单人间 \\
        sport individuel \hspace{6pt} 个人运动项目}\\
    {}&{\hei 近义辨析}&{individuel 与 personnel 近义，individuel 强调"单个的、非集体的"，personnel 强调"私人的、非公共的"}\\
    {}&{\hei 性数变化}&{阴性 individuelle，复数 individuels/individuelles（规则 +s）}\\
    {}&{\hei 位　　置}&{通常置于名词前}\\
\end{vocabEntry}
&
\begin{vocabRemark}
    \vocabMot{individu}{\sposnm}{个人，个体} \\
    \vocabMot{individualisme}{\sposnm}{个人主义} \\
    \vocabMot{individualité}{\sposnf}{个性}
\end{vocabRemark}

\end{listWithRemark}

\newpage
\sectionTitleIV

\begin{listWithRemark}
\frcnTitle
&
{}
\\
\frcn{
    Vite et bon. Et surtout léger. Voilà les mots qui, désormais, sont toujours avec le mot “manger”. Quoique les Français aiment manger, le temps de l’assiette lourde est fini: elle a cédé la place aux plats individuels et faciles à préparer. Les produits naturels sont à l’honneur. Manger sain et léger devient une sorte de bonne action en faveur de la santé. Au bout du compte, les habitudes alimentaires changent vers une nouvelle sorte de joie domestique.
}{
    快捷、美味，尤其是清淡。这些词从此与“吃”形影不离。尽管法国人喜爱美食，但油腻菜肴的时代已经结束：它让位给了个人化且易于准备的菜肴。天然产品备受推崇。吃得健康清淡成为一种有益健康的行为。归根结底，饮食习惯正朝着一种新的家庭乐趣转变。
}
&
\translationRemark[il y a deux jours 表示"两天前"这个过去时间点，不表示从过去持续到现在。]
\\
\cnfrTitle
&
{}
\\
\cnfr{1}{
    她两天前抵达巴黎，为的是在她姑妈莫罗家住两个星期。
}{
    Elle est arrivée à Paris il y a deux jours pour passer deux semaines chez sa tante, la famille Moreau.
}
&
\translationRemark[il y a deux jours 表示"两天前"这个过去时间点，不表示从过去持续到现在。]
\\
\cnfr{2}{
    由于社会发生了变化，人们不再吃油腻的菜肴，而是选择个人份的、易于准备的食物。
}{
    Car la société a changé, on ne mange plus de plats lourds, mais on choisit des plats individuels et faciles à préparer.
}
&
\translationRemark[il y a deux jours 表示"两天前"这个过去时间点，不表示从过去持续到现在。]
\\
\cnfr{3}{
    尽管法国人喜欢去豪华餐厅，但他们也关心保持身材。
}{
    Bien que les Français aiment fréquenter les restaurants de luxe, ils ont aussi le souci de garder leur ligne.
}
&
\translationRemark[il y a deux jours 表示"两天前"这个过去时间点，不表示从过去持续到现在。]
\\
\cnfr{4}{
    她两天前抵达巴黎，为的是在她姑妈莫罗家住两个星期。
}{
    Elle est arrivée à Paris il y a deux jours pour passer deux semaines chez sa tante, la famille Moreau.
}
&
\translationRemark[il y a deux jours 表示"两天前"这个过去时间点，不表示从过去持续到现在。]
\end{listWithRemark}

\end{document}
